\section{Hipóteses do estudo}

Esta pesquisa caracteriza-se como investigação aplicada, multimétodos e orientada à organização e aplicação de um protocolo avaliativo a um sistema RAG em contexto institucional. Nesse contexto, as hipóteses assumem a forma de hipóteses de trabalho verificáveis, funcionando como respostas provisórias ao problema de pesquisa que poderão ser corroboradas, qualificadas ou rejeitadas pelas evidências empíricas produzidas no estudo. Desse modo, não serão tratadas como teste causal clássico entre uma variável independente e uma variável dependente única. Essa opção é compatível com a \textit{Design Science Research}, na qual a construção e a avaliação de artefatos constituem parte do processo de produção de conhecimento \parencite{hevnerEtAl2004designScience,peffersEtAl2007dsrm}.

As hipóteses de trabalho são as seguintes:

\begin{enumerate}[label=\textbf{H\arabic*.}]
  \item \textbf{A pertinência das evidências recuperadas é condição necessária, mas não suficiente, para a verificabilidade das respostas.} A hipótese será examinada pela análise de dois padrões complementares. O primeiro corresponde à ausência de evidência pertinente recuperada para sustentar afirmações presentes na resposta, situação em que a verificabilidade documental fica comprometida. O segundo corresponde aos casos em que há evidência pertinente recuperada, mas a resposta permanece não verificável porque a geração não associa explicitamente as afirmações às evidências recuperadas, utiliza inadequadamente essas evidências ou extrapola o contexto recuperado. A hipótese será considerada corroborada se a análise empírica identificar a ocorrência dos dois padrões em casos documentados. Será qualificada se esses padrões ocorrerem apenas em determinados tipos de pergunta ou sob condições específicas, e será rejeitada se a verificabilidade das respostas puder ser explicada integralmente pela qualidade da recuperação documental.
  \item \textbf{O desempenho do sistema tende a ser inferior em perguntas que exigem integração de múltiplas evidências, raciocínio em etapas ou reconhecimento de informação ausente, em comparação com perguntas factuais simples, delimitadas por autor e apoiadas em evidência documental direta.} A hipótese será examinada por meio de uma bateria de perguntas organizada por categorias operacionais, com comparação de desempenho agregado e análise de erro entre grupos de categorias. Será considerada corroborada se as categorias de maior complexidade apresentarem desempenho inferior ao das categorias simples de forma consistente nos critérios aplicados ao protocolo avaliativo. Será qualificada se essa diferença ocorrer apenas em determinados tipos de pergunta, dimensões avaliativas ou condições específicas, e será rejeitada se o desempenho se mostrar homogêneo entre as categorias ou inverso ao esperado. Essa expectativa direcional deriva de \textit{benchmarks} e estudos de avaliação de RAG, que mostram variação de desempenho conforme a complexidade da pergunta, a necessidade de integração de evidências e a presença de questões não respondíveis \parencite{saadFalconEtAl2024ares,esEtAl2023ragas,zhuEtAl2025rageval,pipitoneAlami2024legalbenchRag}.
  \item \textbf{A avaliação baseada no paradigma \textit{LLM-as-a-judge} apresentará concordância parcial com a avaliação humana, mas será insuficiente como critério único de qualidade.} A hipótese será examinada pela comparação entre julgamentos automatizados e avaliação humana independente. Será considerada corroborada se a avaliação automatizada apresentar concordância parcial com a avaliação humana, mas também revelar divergências suficientes para impedir seu uso como critério único de qualidade. Será qualificada se a concordância variar conforme o tipo de pergunta, a dimensão avaliada ou a formulação da rubrica, e será rejeitada se os julgamentos automatizados se mostrarem plenamente consistentes com a avaliação humana ou, inversamente, incapazes de discriminar diferenças de qualidade segundo os critérios da rubrica. Essa hipótese se fundamenta na literatura sobre avaliação automatizada de linguagem natural e sobre uso de LLMs como avaliadores, que aponta sensibilidade a rubricas, modelos e formato de entrada \parencite{liuEtAl2023geval,liuEtAl2025judgeRag}.

\end{enumerate}

\subsection{Operacionalização e regras de decisão}

A unidade de análise será o par pergunta--resposta, vinculado às evidências recuperadas e às afirmações verificáveis da resposta. Cada pergunta receberá uma categoria de complexidade, um indicador de respondibilidade (respondível ou não respondível) e um conjunto de evidências de referência. As respostas serão avaliadas em escala ordinal de 0 a 2 para correção factual, aderência ao contexto, suporte das citações, completude proporcional e explicitação de insuficiência documental; a pontuação global será a média das dimensões, sem substituir a análise por dimensão.

Para H1, serão formados dois indicadores binários: recuperação pertinente (ao menos uma evidência suficiente em $k=5$) e verificabilidade da resposta (todas as afirmações centrais apoiadas por evidência identificável). A análise separará quatro células: recuperação pertinente/verificável, pertinente/não verificável, não pertinente/verificável e não pertinente/não verificável. H1 será corroborada quando houver pelo menos um caso documentado nas células pertinente/não verificável e não pertinente/não verificável, com exame qualitativo das causas; será qualificada quando os casos se concentrarem em uma categoria ou condição; e será rejeitada quando a célula pertinente/não verificável não ocorrer e a verificabilidade variar apenas com a recuperação.

Para H2, perguntas factuais simples formarão o grupo de referência e perguntas multievidência, multietapas e não respondíveis formarão os grupos de maior complexidade. O contraste primário será a diferença entre as médias de pontuação global e de 	extit{Hit@5} dos grupos. A hipótese será corroborada se o grupo de maior complexidade apresentar desempenho menor em pelo menos duas dimensões e em duas das três comparações principais; será qualificada se a diferença aparecer somente em categorias ou dimensões específicas. Intervalos de confiança por 	extit{bootstrap} e a distribuição dos casos, e não apenas a média, serão reportados.

Para H3, a concordância entre avaliação humana e julgamento automatizado será calculada por $kappa$ ponderado para as dimensões ordinais, percentual de acordo exato e análise dos casos em que o juiz automático aprova uma resposta que os humanos classificam como não verificável. Concordância parcial será descrita como $kappa$ entre 0,40 e 0,75, acompanhada de pelo menos 10\% de divergências em alguma dimensão. H3 será corroborada se houver concordância nesse intervalo e falsos positivos ou falsos negativos documentados em pelo menos 10\% da amostra; será qualificada se a divergência depender da categoria ou da rubrica; e será rejeitada se a concordância for alta e as divergências residuais não alterarem decisões de qualidade.

Em conjunto, essas hipóteses orientam a organização da bateria de perguntas, a seleção dos critérios avaliativos e a interpretação dos resultados empíricos. Seu papel é estruturar a análise da recuperação documental, da geração de respostas e da verificabilidade das fontes, sem antecipar conclusões e sem converter o estudo em um teste causal clássico entre variáveis isoladas.
